\documentclass[11pt]{article}
\usepackage[margin=25mm]{geometry}
\usepackage{fontspec}
\setmainfont{Latin Modern Roman}
\usepackage{amsmath,amssymb,mathtools}
\usepackage{unicode-math}
\setmathfont{Latin Modern Math}
\usepackage{xurl,hyperref}
\hypersetup{colorlinks=true,urlcolor=blue}
\setcounter{secnumdepth}{0}
\setlength{\parindent}{0pt}
\setlength{\parskip}{5pt}
\setlength{\emergencystretch}{3em}
\newcommand{\sref}[2]{\hyperlink{source:#1:#2}{[S#2]}}
\newcommand{\cataloguescope}[1]{\paragraph{Recorded target.}#1}
\newcommand{\R}{\mathbb R}
\newcommand{\C}{\mathbb C}
\newcommand{\N}{\mathbb N}
\newcommand{\Z}{\mathbb Z}
\newcommand{\Q}{\mathbb Q}
\newcommand{\F}{\mathbb F}
\newcommand{\eps}{\varepsilon}
\begin{document}
% UnsolvedMath ID 30006207; source catalogue code OWR-14299088-009
\setcounter{section}{30006206}
\section{Irreducible Components of Landau Varieties}\label{30006207}
{\small UnsolvedMath ID 30006207 · Mathematical Physics}
\subsection{Definitions and mathematical statement}

\paragraph{Shared notation.}$\mathbb N=\{1,2,\ldots\}$; the standard number fields have their usual meanings. All additional conventions are stated locally.


A Landau variety is a locus of kinematic parameters where singularities of an integral, such as a Feynman integral, can arise from degeneracy of its integration geometry. Elimination removes auxiliary integration variables from algebraic criticality equations. An irreducible component is one which cannot be written as the union of two smaller algebraic closed subsets. Feynman-diagram combinatorics includes the graph data defining the Lee--Pomeransky polynomial of the integral.

\paragraph{Question recovered from the primary report.}Determine how to encode the combinatorics of a Feynman diagram in its Landau variety, and whether that encoding assists its irreducible decomposition, as suggested by the Principal Landau Determinants method cited in the report.

\paragraph{Catalogue wording retained for traceability.}
Determine whether the proposed elimination-theoretic construction identifies every irreducible component of the Landau variety, and characterize which candidate components genuinely occur.

\paragraph{Scope and source qualification.}The primary report already proves equality between its elimination construction and the Landau variety (Theorem 3). The catalogue instead recasts the target as completeness and genuine occurrence of candidates. That condensed wording below is retained for traceability; the open primary-source question is the combinatorial encoding and irreducible decomposition stated above.
\subsection{Short English statement}
Encode Feynman-diagram combinatorics in the Landau variety and determine whether it helps decompose the variety into irreducible components.
\subsection{Research attempt}
\paragraph{Model, date and investigated scope.}
OpenAI GPT-6 Astra (configured model \texttt{gpt-6-astra}), ultra reasoning; generated 11 October 2026. Factored a two-propagator stationary discriminant and distinguished algebraic factors from real interior pinches; no general graph-combinatorial decomposition.

\paragraph{Sources and literature check.}
Read the full supplied definition, catalogue wording, source qualifications and listed references for 30006207. No fresh status or novelty search was performed. The argument below proves its auxiliary statements directly and does not import an unverified external theorem.

\paragraph{A one-loop discriminant test.} For positive unequal masses \(m_1,m_2\), consider the two-propagator Feynman-parameter polynomial
\[
 F(x)=xm_1^2+(1-x)m_2^2-sx(1-x)
 =sx^2+(m_1^2-m_2^2-s)x+m_2^2.
\]
An interior stationary zero requires \(F=F'=0\). For \(s\ne0\), elimination of \(x\) gives the quadratic discriminant
\[
 (m_1^2-m_2^2-s)^2-4sm_2^2
 =[s-(m_1+m_2)^2][s-(m_1-m_2)^2].
\]
Thus for fixed distinct positive masses the reduced algebraic candidate locus splits into two distinct linear factors in \(s\).
\paragraph{Real-contour qualification.} The critical point is \(x=(s+m_2^2-m_1^2)/(2s)\). At the plus threshold it equals \(m_2/(m_1+m_2)\in(0,1)\). If \(m_1>m_2\), the minus-threshold point is \(-m_2/(m_1-m_2)<0\); reversing the masses puts it above one. Hence an algebraic component need not meet this real integration interval as an interior pinch. This does not remove it from the algebraic Landau variety or rule out other sheets. Equal masses require a separate \(s=0\) analysis because the displayed division is then invalid.
\paragraph{Corrected source target.} The source already proves its elimination construction equals its Landau variety; that theorem is not reproved or disputed here. The open task is encoding general diagram combinatorics and using it for irreducible decomposition. This single two-edge polynomial gives no graph-to-factor rule across arbitrary diagrams, subgraphs or boundary strata. Establishing such a rule, with coefficient-field and saturation conventions fixed, is the first remaining gap.

\paragraph{Verification and adversarial audit.}
Checked both factors and critical points, unequal positive mass assumptions, real-contour versus algebraic locus and the equal-mass s=0 exception; preserved the source established elimination equality. Model self-audit only; no independent human mathematical verification. Applied contribution prompt A before drafting and B afterwards. Checked both factors and critical points, unequal positive mass assumptions, real-contour versus algebraic locus and the equal-mass s=0 exception; preserved the source established elimination equality. First remaining gap: Develop a diagram-combinatorial rule for irreducible factors across general graphs and boundary strata with the source elimination conventions. The inherited source status is preserved; original-target completion remains zero. Additional GPT-6 Astra reviewer (ultra reasoning): Expanded F(x), eliminated the stationary point for s!=0 and verified the discriminant factors [s-(m1+m2)²][s-(m1-m2)²]. Checked x=m2/(m1+m2) at the upper threshold, x=-m2/(m1-m2)<0 when m1>m2 at the lower threshold, and x>1 on reversing the mass ordering. Confirmed the equal-mass s=0 case is excluded from the division and flagged, and the algebraic component is not incorrectly discarded merely for failing the real-contour test. Compared with the repaired source target: equality of the elimination construction with the Landau variety is treated as supplied known context; graph-combinatorial encoding and irreducible decomposition remain the task.

\paragraph{Tools, process and contributions.}
Python 3 and local file tools used to read definitions and serialize manuscripts. Calculations are symbolic reasoning; no numerical experiment or theorem prover. Contribution prompt A was followed by prompt B. GPT-6 Astra with ultra reasoning drafted the research and performed its stated self-audit; any additional coordinator contribution is identified in the verification and audit. A GPT-6-based Codex coordinator assembled the manuscript and checked repository integration; its exact variant and reasoning configuration were not exposed. No human mathematical contribution, independent human proof check or formal proof certification is supplied. Complete derivations and the consulted source references are included here; private reasoning transcripts are not published.

\paragraph{Final status and first remaining gap.}
\textbf{UNRESOLVED}. Develop a diagram-combinatorial rule for irreducible factors across general graphs and boundary strata with the source elimination conventions.

\section{Completion Estimate}
\noindent\textbf{COMPLETION: 0\%}
This is an elementary auxiliary calculation or an explicit test of a restricted strategy. No new portion of the full source target is established.

\subsection{Sources}
\begin{itemize}
\item[\hypertarget{source:30006207:1}{S1}] ULAM.AI and the UnsolvedMath Contributors, UnsolvedMath, record 30006207 (OWR-14299088-009). Catalogue curation under CC BY 4.0; definitions and English exposition added locally.
\url{https://huggingface.co/datasets/ulamai/UnsolvedMath}.
\item[\hypertarget{source:30006207:2}{S2}] MFO-RIMS Tandem Workshop: Optimization, Theoretical Computer Science and Algebraic Geometry: Convexity and Beyond (2025), Oberwolfach Reports, DOI 10.4171/owr/2025/9, pp. 405--450. Original report and mathematical context. Incorporated report material remains under the source CC BY-SA 4.0 terms unless otherwise noted.
\url{https://doi.org/10.4171/owr/2025/9}.
\end{itemize}
\label{end:30006207}
\end{document}
